\CWHeader{Лабораторная работа \textnumero 4}

\CWProblem{
Необходимо реализовать один из стандартных алгоритмов поиска образцов для указанного алфавита.

{\bfseries Вариант алгоритма:} Поиск одного образца основанный на построении Z-блоков (Z-функция).

{\bfseries Вариант алфавита:} Слова не более 16 знаков латинского алфавита (регистронезависимые).

Запрещается реализовывать алгоритмы на алфавитах меньшей размерности, чем указано в задании.

{\bfseries Формат ввода}

Искомый образец задаётся на первой строке входного файла.
Затем следует текст, состоящий из слов или чисел, в котором нужно найти заданные образцы.
Никаких ограничений на длину строк, равно как и на количество слов или чисел в них, не накладывается.

{\bfseries Формат вывода}

В выходной файл нужно вывести информацию о всех вхождениях искомых образцов в обрабатываемый текст: по одному вхождению на строку.
Следует вывести два числа через запятую: номер строки и номер слова в строке, с которого начинается найденный образец. 
Нумерация начинается с единицы. Номер строки в тексте должен отсчитываться от его реального начала (то есть, без учёта строк, занятых образцами). Порядок следования вхождений образцов несущественен.

{\bfseries Пример}

\begin{tabular}{p{7cm}p{7cm}}
\textbf{Ввод:} & \textbf{Вывод:} \\
\texttt{cat dog cat dog bird} & \texttt{1, 3} \\
\texttt{CAT dog CaT Dog Cat DOG bird CAT} & \texttt{1, 8} \\
\texttt{dog cat dog bird} & \\
\end{tabular}
}
\pagebreak